\scanpage{93}
\deleted{\unsure{suiv.|facteurs}} en groupant suivant les valeurs absolues des zéros et pôles \ldots . Lorsque $X$ est complet et lisse, \inserted{et $F$ \unsure{constant tordu|localement tordu|localement constant}, de poids $p$,} on a $\underline{h}^i(F)=H^{i-p}(\overline X,\overline F)$ comme $\mathbb N$-modules, dans ce cas la connaissance de $L(X,F)$ équivaut à celle des $H^{i\deleted{\unsure{-q|-p|-\rho}}}(\overline X,\overline F)$ et l'opération du frobénius dedans.

\diagram{\begin{tikzcd}[ampersand replacement=\&,column sep=2em,execute at end picture={\draw (current bounding box.south west) -- (current bounding box.north east);\draw (current bounding box.north west) -- (current bounding box.south east);}]
\deleted{X}\arrow[r,dash]\arrow[d,"\deleted{\unsure{f}}"'] \& \deleted{X_{\unsure{\eta|K}}}\arrow[d,dash]\\
\deleted{Y}\arrow[r,dash] \& \deleted{K=\text{corps des fonctions de }Y}
\end{tikzcd}}
\deleted{Soit $f:X\mathbin{-}Y$ morphisme [\unsure{lisse|disons} propre] avec $X,Y$ \unsure{compactifiables|simplifiables} de type fini sur $\mathbb F_p$. Considérons les ouverts denses $U$ de $Y$, et les $L(f^{-1}(U),F)$. Pour $U$ variable, ces fonctions rationnelles}

Considérons d'autre part \inserted{[pour $X$ de dim $n$]} les ouverts \deleted{denses} $U$ de $X$, \inserted{tels que $X-U$ de dim $<n$} et les $L(U,F\mid U)$, qui ne diffèrent \deleted{que} (pour $F$ pur de poids $p$) que par des fonctions rationnelles
\margin{les poids ici \unsure{s'ent|sont|n'est} pour la bonne filtration \deleted{\unsure{ainsi}} cf N\textsuperscript{o} 12}
% slanted
de poids $\leq 2n-2+p$. Donc, à priori la connaissance de l'invariant birationnel $L(X,F)$ mod. facteurs rat. de poids $\leq 2n-2\unsure{+p|{-}p|+\rho}$,
% archive sheet number: 92
